\documentclass[10pt,twocolumn,letterpaper]{article}
\usepackage[review]{cvpr}
\input{preamble}
\definecolor{cvprblue}{rgb}{0.21,0.49,0.74}
\usepackage[pagebackref,breaklinks,colorlinks,allcolors=cvprblue]{hyperref}
\def\paperID{*****}
\def\confName{CVPR}
\def\confYear{2026}
\title{When Agreement Hides Error:\\
Shared Measurement Influence in MRI Slice Interpolation}
\begin{document}
\maketitle

\begin{abstract}
When does prediction agreement cease to be evidence of reliability? Two reconstructions can agree because the same acquired measurements influence them similarly, aligning errors that disappear from their difference. We introduce the Measurement-Influence Reliability Audit, a framework that combines same-target measurement interventions, error-energy decomposition, and local estimator-response analysis without retraining a reconstruction network. A local covariance proposition specifies when shared acquisition variation suppresses disagreement, and motivates an influence statistic that accounts for the strength and direction of responses to physical anchors. Retrospective experiments across three MRI cohorts and five frozen architecture implementations support a reproducible loss of local error observability under sharing; clean shared controls hide approximately 97--99.8\% of paired-error energy in the original two-architecture experiment. Influence refines overlap count where the probe is stable, while fixed-count and architecture failures delimit this interpretation. Case-level ranking is less consistent, and a limited acquired-only scoring experiment does not establish a superior reliability score. The framework exposes the distinction between observable disagreement and hidden common error: agreement alone cannot certify correctness under shared influential evidence.
\end{abstract}

\section{Introduction}
Prediction agreement is often treated as evidence of reliability. It is more informative when alternative predictions fail differently and less informative when they inherit the same error. Our question is: \emph{when does prediction agreement cease to be evidence of reliability?} We study deterministic cross-context disagreement, not all Bayesian, probabilistic, or calibrated uncertainty methods.

MRI slice interpolation offers a controlled setting for this question. A retrospectively available dense examination identifies the missing target plane, the acquired slice positions, and the physical identity of each anchor. Two contexts can reconstruct the same target while sharing or separating acquired measurements. This lets us intervene on evidence dependence rather than merely compare networks with different architectures or accuracy.

\begin{figure*}[t]
\centering\includegraphics[width=.98\textwidth]{p17_qualitative_teaser.pdf}
\caption{Agreement hides spatially aligned error. A median-effect fastMRI case, selected by a recorded rule before new inference, is reconstructed with frozen SwinIR weights. The reference prediction A and target are identical across rows. Shared inputs give similar predictions and small disagreement despite substantial error; disjoint inputs expose more of that error. Every error and disagreement map uses the same scale. These are actual CPU reconstructions, not generated anatomy. The compact chain is conditional, not a universal implication; Proposition 1 states its assumptions. A seeded random case and numerical reproduction differences are in the supplement.}
\label{fig:teaser}
\end{figure*}

Figure~\ref{fig:teaser} shows the failure directly. Under shared anchors, the two predictions smooth the same structures and largely agree on their errors. Disjoint anchors need not make either reconstruction substantially more accurate; they change which error components become visible in the difference. The distinction matters: low disagreement is an observation about two outputs, not a certificate about their unknown target.

The error-difference identity is elementary, and variability missing systematic MRI error is already known~\cite{hu2021predict}. Our contribution is an executable \emph{Measurement-Influence Reliability Audit} (MIRA). It specifies the physical intervention, quantities that require ground truth, local influence computation, stability checks, and unit of statistical inference. A local covariance model explains why shared anchors matter through estimator sensitivity rather than through their count alone. Controlled experiments then test where this explanation succeeds and where it fails.

The distinction between local geometry and global utility is central. A repeatable reduction in observable error does not imply that disjoint disagreement always ranks examinations better. Nor does a stable finite-response probe establish a universal scalar score. We retain the original failures and one additional limited scoring test rather than make successful deployment a prerequisite for a coherent mechanism paper.

Our contributions are: (1) a same-target measurement audit usable with frozen estimators and explicit information-access contracts; (2) a local covariance explanation and an implementation-matched influence measure for common-error blindness; and (3) replication and falsification across cohorts and architecture lineages, separating robust local observability effects from weaker case-ranking consequences. The study is a retrospective synthesis of iteratively developed experiments. Its provenance audit does not establish that every historical choice preceded outcome inspection.

\section{Related Work}
\paragraph{Reconstruction uncertainty and common error.}
Ensembles, dropout, and MRI posterior methods use predictive variation to characterize uncertainty~\cite{lakshminarayanan2017ensembles,gal2016dropout,edupuganti2021uncertainty,narnhofer2022bayesian,luo2023diffusion}. Ensemble diversity interacts with component error~\cite{brown2005diversity,ortega2022diversity}; different architectures need not yield independent failures. Hu \etal explicitly separate random and systematic MRI reconstruction error~\cite{hu2021predict}. We do not claim that common bias is newly discovered. Our distinction is a physical, same-target intervention connecting measurement identity to local response alignment and observable error.

\paragraph{Consistency and interpolation.}
MRI self-consistency can enforce agreement with acquired data~\cite{lustig2010spirit}; it has a different information contract from comparing two missing-slice predictions. SAINT and I3Net learn anisotropic slice synthesis~\cite{peng2020saint,song2024i3net}, while CycleINR represents a different coordinate-based lineage~\cite{fang2024cycleinr}. We audit frozen local implementations, including adapted SAINT-style and internal CycleINR-style models, rather than propose another interpolator. Complementary acquisitions can improve super-resolution~\cite{greenspan2002interslice,nicastro2022shift}, and leave-one-slice-out interpolation checks are established~\cite{horvath2017highorder}. Our question concerns the reliability interpretation of agreement, not novelty of acquisition complementarity or holdout reconstruction.

\section{Our Measurement-Influence Framework}
\subsection{Error observability}
Let $F$ be a frozen reconstruction estimator and $c_a,c_b$ two contexts for target location $z$. On a common support, predictions $f_a,f_b$ have errors $e_a=f_a-y$ and $e_b=f_b-y$. The observable is $D=\|f_a-f_b\|$; the target $y$ is unavailable when forming a deployment score. Retrospectively,
\begin{equation}
D^2=\|e_a\|^2+\|e_b\|^2-2\langle e_a,e_b\rangle .\label{eq:difference}
\end{equation}
Define $T=(\|e_a\|^2+\|e_b\|^2)/2$, common error $\bar e=(e_a+e_b)/2$, and difference error $h=(e_a-e_b)/2$. For $T>0$,
\begin{equation}
S=\|\bar e\|^2/T,\quad O=\|h\|^2/T,\quad S+O=1.\label{eq:energy}
\end{equation}
These are latent evaluation quantities, not estimated probabilities. If $e_a=e_b\ne0$, then $O=D=0$ despite error. Conversely, $\max(\|e_a\|,\|e_b\|)\ge D/2$ gives only a one-sided error bound. Shared fractions and normalized disagreement are algebraically related; their correlation is not independent evidence for a mechanism.

\subsection{Controlled same-target measurement intervention}
The framework takes $F$, target location $z$, two target-preserving contexts, and a registry of physical anchor identities. It records model weights, preprocessing, score support, target geometry, and information used to define each. The intervention holds the target and estimator fixed while varying anchor identity; relevant geometry is matched where feasible and remaining differences are reported.

At $\times4$, phase-0 anchors $\{0,4,8,\ldots\}$ and phase-2 anchors $\{2,6,10,\ldots\}$ are disjoint. Targets with residues 1 or 3 are absent from both. A shared control applies an in-plane flip to the same phase-0 anchors, then maps predictions back. This is a deterministic test-time augmentation (TTA) pair. A disjoint pair instead compares phase-0 and phase-2 reconstructions. It changes measurement dependence together with the companion context; it is not an intervention that fixes every nuisance variable. Nearest-anchor distances can be matched as the multiset $\{1,3\}$ up to reflection, but full signed geometry and context normalization may still differ.

\begin{figure}[t]
\centering
\begin{minipage}{.97\columnwidth}
\hrule\smallskip
\textbf{Algorithm 1: Measurement-Influence Reliability Audit}\par
\smallskip\small
\textbf{Input:} frozen $F$; target locations; anchor-ID registry; declared contexts, preprocessing, support, and perturbation settings. Retrospective targets are a separate evaluation input.\par
\textbf{1. Validate} same-target mapping, acquired IDs, model weights, geometry, and source of masks/scales.\par
\textbf{2. Reconstruct} shared and disjoint pairs; store $f_a,f_b$ and their disagreement on the same support.\par
\textbf{3. Estimate influence} from embedded coefficients or matched anchor responses using Eq.~\ref{eq:influence}.\par
\textbf{4. Check stability} across declared amplitudes/signs; retain failed models as phenomenon-only evidence.\par
\textbf{5. Evaluate retrospectively} alignment, $S,O$, and error ranking; aggregate and resample whole examinations/patients.\par
\textbf{6. Report} all eligible contrasts, uncertainty, failures, and the information-access boundary.\par
\textbf{Output:} an audit of agreement's limitations in this measurement regime; no per-case correctness certificate.\par
\smallskip\hrule
\end{minipage}
\end{figure}

\subsection{Analytical shared influence and a local model}
For $f_c=\sum_j w_{cj}x_j$, embed coefficients in the union of physical anchor IDs, assigning zero to unused anchors. Under independent anchor perturbations with variances $\sigma_j^2$, the shared contribution to output covariance is proportional to $\sum_{j\in c_a\cap c_b}w_{aj}w_{bj}\sigma_j^2$. Overlap count discards sign, magnitude, and noise covariance. Coefficient cosine and common absolute mass $\sum_j\min(|w_{aj}|,|w_{bj}|)$ are descriptive alternatives; common mass alone does not prove positive covariance when coefficients have opposite signs.

\paragraph{Proposition 1 (shared variation under fixed marginals).}
Condition on a fixed target and nominal inputs. Assume the local error model
\begin{equation}
e_a=b_a+J_a\xi_a,\qquad e_b=b_b+J_b\xi_b,\label{eq:local}
\end{equation}
with fixed biases/Jacobians, zero-mean perturbations, marginal covariances $\Sigma_a,\Sigma_b$, and cross-covariance $C=\mathbb E[\xi_a\xi_b^\top]$. Compared with $C=0$ at unchanged marginals and biases, let $\kappa=\operatorname{tr}(J_a C J_b^\top)$. Then
\begin{align}
\Delta\mathbb E\langle e_a,e_b\rangle&=\kappa,\\
\Delta\mathbb E\|f_a-f_b\|^2&=-2\kappa.\label{eq:covariance}
\end{align}
If $\tau=\mathbb E T>0$, the ratio-of-expectations observable fraction $\widetilde O=\mathbb E\|h\|^2/\tau$ changes by $-\kappa/(2\tau)$. Thus positive shared sensitivity increases the expected error inner product and suppresses disagreement energy. Shared variation with opposing sensitivities can do the reverse.

\emph{Proof sketch.} Expanding Eq.~\ref{eq:local}, zero means remove the bias--noise cross terms, leaving $\langle b_a,b_b\rangle+\operatorname{tr}(J_a C J_b^\top)$. Marginal error energies do not depend on $C$; substitute into Eq.~\ref{eq:difference}. The supplement gives the full derivation, remainder qualification, and numerical checks.

This proposition is conditional and local. It does not prove that sharing creates common deterministic bias, that error-vector cosine increases in every realization, or that $\mathbb E O=\widetilde O$. Actual phase interventions can change biases and Jacobians. Their effect therefore requires empirical testing rather than following automatically from the proposition.

\subsection{Neural finite-response shared influence}
For anchor $j$ and a matched zero-mean field $\xi_{js}$, compute $r_{cjs}=F_c(x+\epsilon\xi_{js})-F_c(x)$; absent anchors have zero response. The implemented statistic is
\begin{equation}
\widehat I=
\frac{\sum_s\sum_{j\in c_a\cap c_b}\langle r_{ajs},r_{bjs}\rangle}
{\sqrt{\bigl(\sum_s\sum_{j\in c_a}\|r_{ajs}\|^2\bigr)
\bigl(\sum_s\sum_{j\in c_b}\|r_{bjs}\|^2\bigr)}}.\label{eq:influence}
\end{equation}
It is the cosine of response vectors embedded in physical-anchor space, \emph{not} the average cosine of individual shared-anchor responses. Its denominator includes unshared-anchor energy. Under a local linear response and independent matched anchor fields of covariance $\Sigma_j$, the expected numerator, divided by the number of fields and $\epsilon^2$, is $\sum_j\operatorname{tr}(J_{aj}\Sigma_jJ_{bj}^\top)$, the shared covariance term. The normalized finite-sample ratio is a directional proxy, not an unbiased covariance estimator.

Historical probes use four centered Rademacher fields, amplitude .01, alternatives .005/.02/$-.01$, and an audit threshold of .90 for rank stability. Input/output clipping is part of the implemented finite response. Stability supports numerical reproducibility at those settings, not global Jacobian accuracy. Historical masks and full-examination normalization access target information; Eq.~\ref{eq:influence} can be computed without targets only when its preprocessing and support are also acquired-only. We do not retroactively claim that information contract for the stored probes.

\subsection{Reliability interpretation and scope}
Algorithm 1 requires no estimator retraining. With retrospective targets it can expose how much common error agreement conceals and whether influence explains variation beyond count. Without targets it can measure disagreement, anchor sharing, and suitably defined local response, but it cannot identify an arbitrary common bias. The framework is therefore an audit, not a calibrated detector of every individual failure. Reused test cohorts and outcome-informed analysis choices are explicitly part of its evidence record.

\section{Experimental Design and Provenance}
\paragraph{Cohorts and frozen estimators.}
We use 24 fastMRI brain T2 examinations~\cite{knoll2020fastmri}, 115 IXI brain T2 examinations~\cite{ixi_dataset}, and 22 NingXia knee T2 series from eight patients. Their reconstruction outputs come from existing training splits or patient-group out-of-fold models. fastMRI has one patient shared between historical training and validation, but no test-patient ID overlap; IXI train/validation/test IDs are disjoint. Historical checkpoint-selection bias cannot be quantified from these outputs. Distinct cohorts do not constitute untouched holdouts for this iteratively developed reliability analysis.

HAT/SwinIR supply the three-cohort influence experiment. I3Net, adapted SAINT-style, and internal CycleINR-style add three frozen architecture implementations on both brain cohorts. They are lineage tests, not claims of exact official reproductions. The ten primary checkpoints are hashed in the research ledger. Historical phase-0 wrapper checks reproduce 24 fastMRI examinations per model and one IXI examination per model; they do not verify every IXI wrapper output.

\paragraph{Designs and support.}
The five-architecture contrast uses five missing planes per examination. fastMRI has targets $\{3,5,7,9,11\}$; IXI uses deterministic, approximately residue-balanced targets. These experiments use a GT-derived 1\% mask. Earlier clean HAT/SwinIR controls instead use a phase-0 prediction mask. They are reported separately. A DICOM audit supports slice ordering for fastMRI; its legacy NIfTI affine alone does not establish patient orientation.

The analytical count ladder fixes absolute target distances $\{1,2,3,4\}$ and averages all valid sign configurations at each overlap. The neural pair grid compares reference $(-2,+2)$ with eight contexts, recording span, asymmetry and overlap. Four contexts share exactly one anchor. New four-anchor models use regular $\times4$ grids with overlap zero or three; their support differs from the pair grid. The supplement gives every offset, setting and eligibility rule.

\paragraph{Inference and chronology.}
The independent unit is examination, or patient in NingXia. Plane/context measurements travel with their cluster in resampling. Historical intervals are pointwise unless stated otherwise. Five-model simultaneous intervals use 10,000 joint examination bootstrap draws and a maximum standardized-deviation statistic, separately by cohort and estimand. They do not correct the entire history of hypothesis exploration. Main figures retain null and negative effects.

Our machine-readable audit catalogs source files, attempted analyses, splits, masks, metrics, checkpoints, normalization and bootstrap rules. Repository history does not independently establish all earlier ``prespecified'' assertions. We therefore describe the study as retrospective and remove those assertions. The original ranking-primary architecture protocol, smoke-informed control correction, and subsequent emphasis on local observability remain documented. Qualitative case selection and the one new score family were recorded before their new evaluation, but neither creates a new independent cohort.

\section{Results}
\subsection{Phenomenon: common error becomes invisible}
Figure~\ref{fig:control} and Table~\ref{tab:controlled} summarize the original two-architecture clean comparison. Mean observable fractions under sharing are only 2.61/1.58/.25\% on fastMRI/IXI/NingXia. Pooled-plane error-ranking contrasts remain positive with clustered intervals. These are not examination-ranking estimates: pooled planes contribute to each point correlation, while patients/examinations determine interval resampling. The five-architecture case-ranking test is reported separately below.

\input{p17_control_table.tex}
\begin{figure*}[t]
\centering\includegraphics[width=.95\textwidth]{fig2_controlled_intervention.pdf}
\caption{Controlled sharing and error observability. Left: pooled-plane disagreement ranking against fixed phase-0 error, with examination/patient-clustered pointwise 95\% intervals. Right: mean shared and observable paired-error fractions in clean shared HAT/SwinIR controls. Fractions describe error geometry; they are not independent evidence for the algebraic identity.}
\label{fig:control}
\end{figure*}

The objectively selected example in Fig.~\ref{fig:teaser} makes the spatial failure visible. Shared prediction MAEs are 73.8 and 77.7 but their difference MAE is 22.8; observable fraction is 2.0\%. With the same target and reference prediction, disjoint difference MAE rises to 69.3 and observable fraction to 19.8\%, while paired error changes little. These values describe that example, not an estimate of prevalence. Its CPU predictions differ slightly from historical mixed-precision outputs; the cohort statistics are unchanged.

\subsection{Mechanism: influence refines count conditionally}
Raw overlap has an overall trend but fails as a monotone dose: the compact analytical ladder reverses from zero to one shared anchor in every cohort. Unequal coefficient influence improves alignment association, explaining why two contexts with the same count can differ. This rejects count sufficiency, not every possible dependence mechanism.

Figure~\ref{fig:influence} combines broad-grid association with the stricter fixed-count test. Source review found that the historical broad grid included reference self-pairs, which mechanically have unit influence and error cosine. A post-hoc integrity sensitivity removes all 966 self-pairs: influence still improves absolute error-alignment correlation over count by .060--.128 across six strata, with six-stratum simultaneous clustered intervals above zero. At exactly one shared anchor, adjusted influence associations remain positive on both brain cohorts but are null/mixed on NingXia (Table~\ref{tab:neural}). The eight-patient failure is not relabeled as merely lower precision.

\begin{figure*}[t]
\centering\includegraphics[width=.98\textwidth]{p17_influence_beyond_count.pdf}
\caption{Beyond-count evidence and its boundary. (a) Broad-grid advantage after removing self-pairs, with six-stratum simultaneous 95\% intervals. (b) Influence--alignment partial correlation at exactly one shared anchor, controlling span, asymmetry and examination; intervals are historical pointwise clustered 95\%. The NingXia fixed-count result remains null/mixed. Original self-inclusive summaries and other endpoints are retained in the supplement.}
\label{fig:influence}
\end{figure*}
\input{p17_influence_table.tex}

The HAT/SwinIR stability audit retains 42 passing numerical settings (minimum rank correlation .942), but only one of three physical-plane pair gates passes. New-architecture SAINT passes amplitude/sign stability (minimum .940), yet its alignment advantage over count is .014 [$-.002$,.029] and fixed-count estimate .063 [$-.021$,.149]. I3Net and CycleINR fail the .90 stability gate. They support the phenomenon, not a stable finite-response influence claim. Numerical stability and beyond-count evidence are distinct requirements.

\subsection{Replication: local geometry is stronger than ranking}
Across five architectures, fastMRI local observable gains range from 8.0 to 22.6 percentage points and IXI gains from 25.2 to 30.2. All five-model simultaneous intervals remain positive, including all four within-/between-case mean/median sensitivity combinations. Every IXI case effect is positive for each model. Reflection-balanced identity-swap estimates remain positive, while all five mirror-difference intervals include zero; this constrains a particular geometry alternative, not all confounding.

\begin{figure*}[t]
\centering\includegraphics[width=.96\textwidth]{p13_architecture_robustness.pdf}
\caption{Local replication versus case-level ranking. Local effects retain positive five-model simultaneous intervals; fastMRI local points use means and IXI medians. Ranking uses within-case medians in both cohorts. Thin intervals are simultaneous and thick intervals historical pointwise intervals. Only three IXI ranking intervals and no fastMRI ranking intervals exclude zero simultaneously. Models and cohorts are not independent replicate experiments to be pooled.}
\label{fig:arch}
\end{figure*}

The weaker downstream result is examination ordering. Under simultaneous inference, ranking gains exclude zero for HAT, SwinIR and SAINT on IXI, but not I3Net/CycleINR and not any fastMRI model. A local mechanism can be reproducible while a global ranking law fails. The historical shared-threshold stable-wrong rule identifies only four of 120 shared CycleINR fastMRI planes and none for the other models; those thresholds came from the reported test distribution, not external detector validation.

\paragraph{Practical uncertainty constructions.}
The shared neural control itself is a two-view TTA consistency construction. A new, acquired-mask experiment additionally tests a two-member Linear/Cubic ensemble on 115 IXI examinations. Changing only the companion phase improves fixed-reference-error ranking for ensemble variance by .031 [.012,.050] and .036 [.016,.056] under two-direction simultaneous intervals; local gains are about 32.1 points. Disagreement MAE is reported alongside it. Population variance equals the squared pair difference divided by four, so these are related proxies, not two independent discoveries. The same complete-support protocol is infeasible in shallow fastMRI volumes and they are excluded by geometry. This extends the measurement test to a practical analytic ensemble, not to all neural ensembles or calibrated uncertainty methods.

\subsection{A bounded attempt at actionable scoring}
We test one acquired-only guard against small-disagreement failures: $G=\max\{F_{\rm dev}(D_n),F_{\rm dev}(A_n)\}$, where $A$ is Anchor-CV, $n$ denotes acquired-range normalization, and each empirical CDF is fitted on 12 development fastMRI cases. It adds an acquired-data floor when predictors reuse phase-0 anchors; it is not a learned influence score. Two fixed model pairs cover HAT/SwinIR and I3Net/CycleINR. Thresholds and the single formula are then applied to the other 12 fastMRI cases and all 115 IXI cases, compared with raw normalized disagreement, adjacent MAE/RMSE, Anchor-CV, range and spacing.

Neither result supports a new method claim. The fixed both-high-error/low-disagreement label has no positive internal-evaluation cases and only two per pair on IXI. Internal AUC/AP are therefore undefined, and all 12 IXI guard-minus-baseline simultaneous AP intervals include zero. Historical cohort reuse further prevents an untouched-holdout interpretation. We stop after this candidate. The mechanism is reproducible within its scope, while a universally superior scalar reliability score remains open. Full values, degenerate resamples and earlier failed candidates are retained in the supplement.

\section{Discussion and Limitations}
MIRA separates three evidence levels. \emph{Replicated locally:} shared measurement regimes can align errors and conceal common error energy. \emph{Context dependent:} stable influence contains information beyond count on some cohorts and architectures, and disjoint contexts sometimes improve case ranking. \emph{Unestablished:} calibrated correctness, universal scalar superiority, prospective acquisition benefit, or non-MRI generalization. The covariance proposition explains a conditional sensitivity mechanism; it neither establishes global causality nor turns a retrospective audit into a target-free certificate.

The central limitations are concrete. Analyses repeatedly inspect the same cohorts, some controls and endpoints evolved after results, and filesystem dates do not establish prespecification. Historical influence masks and normalization use targets. Interventions leave residual geometry, boundary and preprocessing differences. NingXia has eight independent patients; the five-architecture study covers brain T2, not a matched neural knee replication. Adapted implementations and imperfect historical training/validation separation limit architecture interpretation. Statistical intervals quantify sampling uncertainty conditional on the chosen analyses, not selection uncertainty from the full research history.

The complete 13-model reconstruction benchmark, training-budget audit, mechanism-chain associations, retrospective densification, normalization variants, and external analytical MRA studies remain in the supplement. They contextualize scope rather than establish a universal method. In particular, current-error adjustment weakens Anchor-CV benefit associations, Res-SRDiff retains an adverse boundary, and anatomy-specific ranking need not improve general error ranking. No clinical, prospective acquisition, or natural-image generalization claim follows from this work.

\section{Conclusion}
Agreement between frozen reconstructions can conceal common error inherited from shared influential measurements. A same-target audit, a conditional covariance model, and replication with explicit failures make that limitation measurable. Local observability is more reproducible than examination ranking, and practical target-free superiority remains unproven.

\clearpage
{\small\bibliographystyle{ieeenat_fullname}\bibliography{references}}
\end{document}
